\subsection{Every act cone over a monoid of order at most nine}

\begin{theorem}\label{C:nine-act-cones}
For every finite monoid $M$ with $|M|\leq9$ and every nonempty
finite right $M$-set $S$, the category $\mathcal C(M,S)$ does not
have classifying space of type $\mathbb{CP}^2$.  There is no bound
on the cardinality of $S$.  The same exclusion holds for triangular
biset categories whose other endomorphism monoid has projective
trivial rational module on the side used for coinvariants.
\end{theorem}

\begin{proof}
Orders at most eight are covered by
Proposition~\ref{C:eight-act-cones}.  For order nine, if
$H^2(BM;\mathbb Q)=0$, every degree-two class of the act cone lifts
relative to $BM$.  The quotient by $BM$ is a suspension, so all
such products vanish, as in
Proposition~\ref{C:seven-act-cones}.  We may therefore assume this
vertex cohomology is nonzero.

Write the minimal ideal in Rees form with $r$ rows, $s$ columns,
and maximal group $H$.  If $r=1$ or $s=1$, the monoid has finite-group
classifying-space type and its positive rational cohomology is zero.
If $r\leq2$, Theorem~\ref{C:two-row-act-cups} applies directly.
The remaining size inequality is $rs|H|+1\leq9$ with $r\geq3$
and $s\geq2$.  It forces $H=1$, and leaves the shapes $3\times2$
or $4\times2$.  In the latter case the monoid is exactly the
rectangular band with an identity adjoined.  If $K$ is its band,
the ideal $\mathbb QK$ is projective on the left, and
$\mathbb QM/\mathbb QK=\mathbb Q$.  The row-augmentation kernel
in the projective cover $\mathbb QM e\to\mathbb Q$ is a direct
sum of copies of this quotient.  Thus
$\operatorname{pd}_{\mathbb QM}\mathbb Q\leq2$; the act-cone Tor
formula excludes degree-four homology.

For shape $3\times2$, the complete translation-table enumeration
has $191$ oriented order-nine tables, including reducible targets.
The exact, independently verified partition is
\[
\begin{array}{l|r}
 \text{all-act obstruction}&\text{number of tables}\\\hline
 H^2(BM;\mathbb Q)=0&62\\
 \operatorname{pd}_{\mathbb QM}\mathbb Q\leq2&73\\
 \text{exact periodic projective tail}&23\\
 \Omega^2\mathbb Q\text{ projective}&4\\
 \text{stable retract into a later syzygy}&22\\
 \text{split rational row augmentation}&7
\end{array}
\]
These counts sum to $191$.  The first two lines give the obstructions
already described.  The periodic and stable-retract lines force
nonzero degree-four act-cone homology to recur in higher degrees;
use the exact tensor formula of Lemma~\ref{C:act-cone-Tor} and
Lemma~\ref{C:stable-retract-Tor}.  A projective second syzygy is
another projective-dimension-two certificate.  Every module action,
projective differential, exact kernel, periodicity identification,
and stable identity is checked in the standalone verifier with
summary \path{code/outputs/b_nine_small_acts_verified.json}; the
source records retain the literal tables and exact certificate
locations.  The seven final tables are precisely those in
Proposition~\ref{C:nine-three-row-all-acts}, which excludes all their
finite acts by the split-character argument.  No transfer of an
arbitrary act along a mere homotopy equivalence of monoids is used.

For the biset assertion, take the actual finite orbit act of the
other monoid.  Exact rational coinvariants make the quotient functor
a rational homology equivalence on the edge bar and on the second
vertex, hence on the complete homotopy pushout.  It induces an
isomorphism of rational cohomology rings.  This is the same
comparison as Corollary~\ref{C:two-row-biset-cups}, now applying
the full order-nine act-cone exclusion to the quotient.
\end{proof}
